\section{补充结果与复现说明}\label{app:results}
表~\ref{tab:service} 补充主缓冲扫描的服务分项。冷启动包括证书生成及检查、编码或计划和编译。缺失计算为未命中完整答案缓存时的计算，输出列为形成完整支持列表的记录成本；各分项独立取中位数，且存在额外管理开销，所以不能将列相加作为总计。$C=2$ 有三次全局缓存命中，其余主点均无命中。

\begin{table}[p]
\centering\small
\caption{主缓冲扫描的服务成本细项，单位秒。$r/C$ 为轮数／容量，$U$ 为不同假设数，答案为不同完整支持答案数。成功点三次中位数；未完成不填最终请求统计。}\label{tab:service}
% Generated from frozen JSON; do not edit.
\begin{tabular}{llrrrrrr}
\toprule
$r/C$ & 方法 & $U$ & 答案 & 冷启动 & 缺失计算 & 输出 & 总计 \\
\midrule
4/3 & B0 & 1000 & 726 & 0.0462 & 0.9947 & 0.0043 & 1.0472 \\
4/3 & B1 & 1000 & 726 & 0.0218 & 0.7169 & 0.0037 & 0.7481 \\
4/3 & BD & 1000 & 726 & 0.0272 & 0.6446 & 0.0038 & 0.6771 \\
4/3 & B2 & 1000 & 726 & 0.0373 & 0.1769 & 0.0035 & 0.2223 \\
8/3 & B0 & 1000 & 991 & 0.3900 & 16.3947 & 0.0081 & 16.7957 \\
8/3 & B1 & 1000 & 991 & 0.0745 & 12.0800 & 0.0079 & 12.1646 \\
8/3 & BD & 1000 & 991 & 0.1009 & 11.7976 & 0.0088 & 11.8927 \\
8/3 & B2 & 1000 & 991 & 0.1238 & 0.5729 & 0.0059 & 0.7028 \\
16/3 & B0 & -- & -- & -- & -- & -- & 超时 \\
16/3 & B1 & -- & -- & -- & -- & -- & 超时 \\
16/3 & BD & -- & -- & -- & -- & -- & 超时 \\
16/3 & B2 & 1000 & 991 & 0.3290 & 1.5939 & 0.0197 & 1.9647 \\
32/3 & B0 & -- & -- & -- & -- & -- & 超时 \\
32/3 & B1 & -- & -- & -- & -- & -- & 超时 \\
32/3 & BD & -- & -- & -- & -- & -- & 超时 \\
32/3 & B2 & 1000 & 991 & 0.8506 & 4.7251 & 0.0885 & 5.6706 \\
8/2 & B0 & 997 & 991 & 0.3082 & 6.8212 & 0.0080 & 7.1388 \\
8/2 & B1 & 997 & 991 & 0.0633 & 3.8329 & 0.0069 & 3.9083 \\
8/2 & BD & 997 & 991 & 0.0890 & 3.5270 & 0.0076 & 3.6265 \\
8/2 & B2 & 997 & 991 & 0.0373 & 0.2393 & 0.0061 & 0.2857 \\
8/4 & B0 & 1000 & 991 & 0.3377 & 28.2518 & 0.0090 & 28.6004 \\
8/4 & B1 & 1000 & 991 & 0.0842 & 20.9540 & 0.0085 & 21.0500 \\
8/4 & BD & 1000 & 991 & 0.1028 & 22.7015 & 0.0100 & 22.8425 \\
8/4 & B2 & 1000 & 991 & 1.2137 & 2.9863 & 0.0090 & 4.2130 \\
8/6 & B0 & 1000 & 991 & 0.3735 & 47.0278 & 0.0094 & 47.4084 \\
8/6 & B1 & 1000 & 991 & 0.1039 & 37.3728 & 0.0091 & 37.4896 \\
8/6 & BD & 1000 & 991 & 0.1324 & 38.1007 & 0.0097 & 38.2462 \\
8/6 & B2 & -- & -- & -- & -- & -- & 状态上限 \\
\bottomrule
\end{tabular}

\end{table}

\begin{table}[t]
\centering\footnotesize
\caption{$16$ 轮机制对照。时间为秒，查询部分为 $100$ 次完整支持请求。前三行来自冻结的机制批次，最后两行来自匹配的长写间隔对照，各行均三次中位数；只在对应组内比较。}\label{tab:mechanism}
\resizebox{\linewidth}{!}{% Generated from frozen JSON; do not edit.
\begin{tabular}{lrrrrrrrrr}
\toprule
表示／布局 & $K$ & 组 & $S_{\max}$ & 状态 & 转移 & $T_{\rm plan}$ & $T_{\rm compile}$ & $T_{\rm query}$ & 总计 \\
\midrule
稀疏／专用 & 13 & 913 & 10 & 1845 & 1970 & 0.1076 & 0.0781 & 0.1319 & 0.5164 \\
稀疏／普通 & 15 & 913 & 10 & 1877 & 2002 & 0.1073 & 0.0779 & 0.1349 & 0.4935 \\
稠密／专用 & 162 & 3045 & 10 & 5117 & 5242 & 0.8139 & 1.5881 & 0.2014 & 2.8777 \\
长间隔／专用 & 15 & 965 & 10 & 1946 & 2071 & 0.0955 & 0.0801 & 0.1348 & 0.5250 \\
长间隔／普通 & 17 & 965 & 10 & 1978 & 2103 & 0.0923 & 0.0897 & 0.1281 & 0.4889 \\
\bottomrule
\end{tabular}
}
\end{table}

输入、实现与原始测量保存在仓库 \url{https://github.com/ChenKaichen-SCUT/Work}。论文源包附有冻结数据快照；\texttt{generate\_results.py} 只读取这些 JSON，自动生成四张实验图及所有结果表。材料版本、输入摘要和分批环境记录见源包中的独立复现说明。

用 Python、Matplotlib 和 NumPy 重建图表，然后用 XeLaTeX、BibTeX 编译：
\begin{verbatim}
python generate_results.py --snapshot
xelatex -interaction=nonstopmode -halt-on-error main.tex
bibtex main
xelatex -interaction=nonstopmode -halt-on-error main.tex
xelatex -interaction=nonstopmode -halt-on-error main.tex
\end{verbatim}
源文件使用 TeX Live 所带 Fandol 中文字体，图中使用 DejaVu Sans。各实验批次保留各自的输入、源码摘要和测量口径；重绘图表不改变实验版本。完整源码已包含生成的矢量图和表，直接编译论文无须安装实验求解器。
